\subsection{DEFINITION OF THE TENSOR ALGEBRA OF A MODULE}

Let A be a commutative ring and M an A-module. For every integer \(n \geqslant 0\) , the A-module the tensor product of n modules equal to M (also called the n-th tensor power of M) is denoted by \(\bigotimes^{n} M\) , or \(M^{\otimes n}\) , or \(T^{n}(M)\) , or \(T_{A}^{n}(M)\) , or \(\operatorname{Tens}^{n}(M)\) ; then \(T^{1}(M) = M\) ; also we write \(T^{0}(M) = A\) . The A-module the direct sum \(\bigoplus_{n \geqslant 0} T^{n}(M)\) is denoted by \(T(M)\) or \(\operatorname{Tens}(M)\) . We shall define a graded A-algebra structure of type N on \(T(M)\) , by defining for every ordered pair of integers \(p \geqslant 0\) , \(q \geqslant 0\) , an A-linear mapping

\[
m _ {p q} \colon \mathsf {T} ^ {p} (\mathbf {M}) \otimes_ {\mathbf {A}} \mathsf {T} ^ {q} (\mathbf {M}) \rightarrow \mathsf {T} ^ {p + q} (\mathbf {M})
\]

(§ 3, no. 1, Remark). For \(p > 0\) and \(q > 0\) , \(m_{pq}\) is the associativity isomorphism (II, § 3, no. 9) and, when \(p = 0\) (resp. \(q = 0\) ), \(m_{0,q}\) is the canonical isomorphism of \(A \otimes_A T^q(M)\) onto \(T^q(M)\) (resp. \(m_{p,0}\) is the canonical isomorphism of \(T^p(M) \otimes_A A\) onto \(T^p(M)\)) (II, § 3, no. 4, Proposition 4). Then, for \(x_t \in M\) , \(\alpha \in A\) ,

\[
\left\{ \begin{array}{r l} (x _ {1} \otimes \dots \otimes x _ {p}) \cdot (x _ {p + 1} \otimes \dots \otimes x _ {p + q}) & \\ = x _ {1} \otimes \dots \otimes x _ {p} \otimes x _ {p + 1} \otimes \dots \otimes x _ {p + q} \\ \alpha \cdot (x _ {1} \otimes \dots \otimes x _ {p}) = \alpha (x _ {1} \otimes \dots \otimes x _ {p}). \end{array} \right.\tag{1}
\]

It is immediate that the multiplication thus defined on \(\mathsf{T}(\mathbf{M})\) is associative and admits as unit element the unit element \(1\) of \(\mathbf{A} = \mathsf{T}^0 (\mathbf{M})\) .

DEFINITION 1. For every module M over a commutative ring A, the tensor algebra of M, denoted by T(M), or Tens(M), or \(\mathsf{T}_{\mathbf{A}}(\mathbf{M})\) , is the algebra \(\bigoplus_{n\geqslant 0}\mathsf{T}^n (\mathbf{M})\) with the multiplication defined in (1). The canonical injection \(\phi :\mathsf{T}^{1}(\mathbf{M})\to \mathsf{T}(\mathbf{M})\) (II, § 1, no. 12) (also denoted by \(\phi_{\mathbf{M}}\) ) is called the canonical injection of M into T(M).

PROPOSITION 1. Let E be a (unital) A-algebra and \(f:M\to E\) an A-linear mapping. There exists one and only one A-algebra homomorphism \(g:T(M)\to E\) such that \(f=g\circ\phi\) .

In other words, \((\mathsf{T}(\mathbf{M}), \phi)\) is a solution of the universal mapping problem (Set Theory, IV, §3, no. 1), where \(\Sigma\) is the species of A-algebra structure, the \(\alpha\) -mappings being the A-linear mappings from the module M to an A-algebra. Observe that here there is no question of a graduation on \(\mathsf{T}(\mathbf{M})\) .

For every finite family \((x_{i})_{1\leqslant i\leqslant n}\) of n elements of M, by definition of the product in T(M), \(x_{1}\otimes x_{2}\otimes\cdots\otimes x_{n}=\phi(x_{1})\phi(x_{2})\ldots\phi(x_{n})\) ; then necessarily \(g(x_{1}\otimes x_{2}\otimes\cdots\otimes x_{n})=f(x_{1})f(x_{2})\ldots f(x_{n})\) for \(n\geqslant1\) and \(g(\alpha)=\alpha e\) (if e is the unit element of E) for \(\alpha\in A\) , which proves the uniqueness of g. Conversely, note that, for all n\textgreater0, the mapping

\[
(x _ {1}, \dots , x _ {n}) \mapsto f (x _ {1}) f (x _ {2}) \dots f (x _ {n})
\]

of \(\mathbf{M}^n\) into \(\mathbf{E}\) is A-multilinear; hence there corresponds to it an A-linear mapping \(g_n: \mathsf{T}^n(\mathbf{M}) \to \mathbf{E}\) such that

\[
g \left(x _ {1} \otimes x _ {2} \otimes \dots \otimes x _ {n}\right) = f \left(x _ {1}\right) f \left(x _ {2}\right) \dots f \left(x _ {n}\right).\tag{2}
\]

Also we define the mapping \(g_0 \colon \mathsf{T}^0(\mathbf{M}) \to \mathbf{E}\) as equal to \(\eta_{\mathbf{E}}\) (§ 1, no. 3), in other words \(g_0(\alpha) = \alpha e\) for \(\alpha \in \mathbf{A}\) . Let \(g\) be the unique A-linear mapping of \(\mathsf{T}(\mathbf{M})\) into \(\mathbf{E}\) whose restriction to \(\mathsf{T}^n(\mathbf{M})\) is \(g_n\) ( \(n \geqslant 0\) ); it is immediate that \(g \circ \phi = g_1 = f\) and it remains to verify that \(g\) is an A-algebra homomorphism. By construction \(g(1) = e\) and it suffices by linearity to show that \(g(uv) = g(u)g(v)\) for \(u \in \mathsf{T}^p(\mathbf{M})\) and \(v \in \mathsf{T}^q(\mathbf{M})\) ( \(p > 0, q > 0\) ); now it follows from formulae (1) and (2) that this relation is true when \(u = x_1 \otimes x_2 \otimes \cdots \otimes x_p\) and \(v = x_{p+1} \otimes \cdots \otimes x_{p+q}\) (where the \(x_i\) belong to \(\mathbf{M}\) ). It is therefore true for \(u \in \mathsf{T}^p(\mathbf{M})\) and \(v \in \mathsf{T}^q(\mathbf{M})\) by linearity.

Remark. Suppose that \(\mathbf{E}\) is a graded A-algebra of type \(\mathbf{Z}\) , with graduation \((\mathbf{E}_n)\) , and suppose also that

\[
f (\mathbf {M}) \subset \mathrm{E} _ {1}.\tag{3}
\]

Then it follows from (2) that \(g(\mathsf{T}^p (\mathbf{M})) \subset \mathrm{E}_p\) for all \(p \geqslant 0\) and hence \(g\) is a graded algebra homomorphism.

